%==============================================================================
\section{Application: Hypothetical Thrombolytic MRCT Using GUSTO-I Data}\label{sec:application}
%==============================================================================

We illustrate the nABCD framework through a hypothetical scenario grounded in publicly available data from GUSTO-I, a large international trial in acute myocardial infarction (AMI). A sponsor is developing a novel thrombolytic agent (Drug~T) for AMI and is planning a Phase~3 MRCT across multiple regions. At the planning stage, the sponsor uses IPD from GUSTO-I ($N = 40{,}830$; 16 regions)\cite{gusto1993} to assess distributional similarity of candidate effect modifiers across regions and to identify suitable pooling partners for a designated anchor region. GUSTO-I was not designed as an MRCT and its regions are anonymized; it serves as a convenient publicly available data source for methodological demonstration, and the anonymization precludes demographic or geographic interpretation of the distributional patterns. In practice, sponsors would assemble regional effect modifier distributions from the most representative and current sources available.

\subsection{Scenario Design and Effect Modifier Selection}\label{sec:app_scenario}

For this illustrative exercise, the sponsor identifies two candidate effect modifiers for Drug~T's treatment effect in AMI: age and systolic blood pressure (SBP). Drug~T is a novel agent. Its Phase~2 programme has suggested age and SBP as candidate effect modifiers on the basis of exploratory analyses and biological plausibility; the precision of these analyses is insufficient for numerical estimation of $L$. At this planning stage the sponsor therefore has no quantitative prior on the CATE sensitivity $L$ for either candidate effect modifier, and class-level indirect evidence is of limited value as a numerical substitute. For \textbf{age}, the Fibrinolytic Therapy Trialists' (FTT) collaborative meta-analysis of existing fibrinolytic agents\cite{ftt1994} concluded benefit is present \emph{``irrespective of age''} and does not report a per-year CATE slope on the absolute mortality scale, so class transferability would not yield a defensible numerical prior for Drug~T. For \textbf{SBP}, no class-level quantitative estimate of CATE sensitivity is available at all. We therefore treat $L$ as unknown a priori for both candidate effect modifiers, and apply the $L^*$ reverse-calculation pathway (equation~\ref{eq:lstar}) to both.

The sponsor selects Region~8 as the anchor region---representing, for instance, a regulatory market where local efficacy evidence is required but sample size is limited. Region~8 comprises $n = 2{,}916$ patients in GUSTO-I, with baseline characteristics summarized in Table~\ref{tab:gusto_r8}. The exercise of identifying Region~8's suitable pooling partners simulates the decision a sponsor would face when seeking regional efficacy assessment through pooling as described in ICH E17.\cite{iche17}

\begin{table}[ht]
\caption{Region~8 baseline characteristics (GUSTO-I).\label{tab:gusto_r8}}
\begin{tabular*}{\textwidth}{@{\extracolsep\fill}lcccc@{}}
\toprule
\textbf{Candidate effect modifier} & \textbf{Mean} & \textbf{SD} & \textbf{Skewness} & \textbf{IQR\textsubscript{pooled}} \\
\midrule
Age (years) & 60.2 & 12.1 & $-0.17$ & 17--18 \\
SBP (mmHg) & 132.4 & 22.9 & 0.20 & 30--34 \\
\bottomrule
\end{tabular*}
\begin{tablenotes}
\item $n = 2{,}916$. IQR\textsubscript{pooled} ranges reflect variation across the 15 partner-region pairs.
\end{tablenotes}
\end{table}

\subsection{Distributional Assessment: nABCD Across 15 Partner Regions}\label{sec:app_nabcd}

We computed nABCD between Region~8 and each of the remaining 15 regions for both candidate effect modifiers, with percentile bootstrap 95\% confidence intervals ($B = 2{,}000$). Table~\ref{tab:gusto_nabcd} presents the results. Figure~\ref{fig:gusto_forest} displays these results as a forest plot with bootstrap CIs, with partners ordered by ascending age nABCD to aid visual comparison.

\begin{table}[ht]
\caption{nABCD values (with 95\% percentile bootstrap CIs) for Region~8 versus each partner region (GUSTO-I), listed by region number.\label{tab:gusto_nabcd}}
\begin{tabular*}{\textwidth}{@{\extracolsep\fill}lrll@{}}
\toprule
\textbf{Partner} & $n$ & \textbf{nABCD\textsubscript{age} [95\% CI]} & \textbf{nABCD\textsubscript{SBP} [95\% CI]} \\
\midrule
R1  & 2188 & 0.018 [0.014, 0.028] & 0.058 [0.044, 0.079] \\
R2  & 2952 & 0.061 [0.044, 0.079] & 0.015 [0.012, 0.030] \\
R3  & 2030 & 0.076 [0.057, 0.095] & 0.051 [0.035, 0.074] \\
R4  & 2876 & 0.016 [0.010, 0.032] & 0.042 [0.029, 0.060] \\
R5  & 1909 & 0.011 [0.010, 0.026] & 0.056 [0.038, 0.077] \\
R6  & 1585 & 0.026 [0.018, 0.040] & 0.050 [0.032, 0.076] \\
R7  & 3150 & 0.011 [0.008, 0.026] & 0.082 [0.064, 0.098] \\
R9  & 3123 & 0.017 [0.011, 0.032] & 0.110 [0.091, 0.130] \\
R10 & 1717 & 0.038 [0.022, 0.058] & 0.062 [0.045, 0.085] \\
R11 & 2491 & 0.035 [0.021, 0.052] & 0.104 [0.082, 0.122] \\
R12 & 4352 & 0.024 [0.013, 0.040] & 0.097 [0.083, 0.117] \\
R13 & 2297 & 0.034 [0.022, 0.051] & 0.037 [0.023, 0.061] \\
R14 & 3437 & 0.020 [0.012, 0.036] & 0.064 [0.052, 0.082] \\
R15 & 2576 & 0.017 [0.011, 0.034] & 0.074 [0.057, 0.093] \\
R16 & 1231 & 0.050 [0.034, 0.071] & 0.071 [0.051, 0.095] \\
\bottomrule
\end{tabular*}
\begin{tablenotes}
\item Point estimates with 95\% percentile bootstrap CIs ($B = 2{,}000$, seed = 2026). Regions are listed by region number; ordering carries no implicit ranking. Bootstrap CIs quantify estimation uncertainty; narrow CIs lend greater confidence to individual nABCD values, while overlapping CIs indicate partners whose relative position is not sharply resolved by the data.
\end{tablenotes}
\end{table}

\begin{figure}[ht]
\centering
\includegraphics[width=\textwidth]{figures/fig_gusto_r8_forest}
\caption{Forest plot of nABCD point estimates with 95\% percentile bootstrap CIs for Region~8 versus each partner region, for age (left) and systolic blood pressure (right). Partners are ordered by ascending age nABCD to aid visual comparison; see Table~\ref{tab:gusto_nabcd} for values listed by region number. The bootstrap CI widths indicate the precision of each partner's estimated nABCD.}
\label{fig:gusto_forest}
\end{figure}

The two candidate effect modifiers produce different patterns of distributional similarity. For age, partner-region nABCD point estimates span a narrow range, from 0.011 (R5, R7) at the low end up to 0.076 (R3) at the high end; eleven of the fifteen partners have age nABCD below 0.040. For SBP, point estimates span a wider range, from 0.015 (R2) up to 0.110 (R9), with the majority of partners falling in the 0.050--0.110 interval. SBP distributions therefore exhibit greater between-region heterogeneity than age distributions across GUSTO-I, and the ranking of partners by each candidate effect modifier differs accordingly.

Several partners fall in the middle of each ranking where uncertainty intervals overlap substantially. R16's age nABCD ($0.050$) is the 13th-smallest value with 95\% CI $[0.034, 0.071]$, placing it in a region where its relative position is not sharply determined by the data---the CI spans the range covered by several partners with smaller point estimates. R6's age nABCD ($0.026$) is the 9th-smallest value but its SBP nABCD of $0.050$ sits near the lower end of the SBP-nABCD distribution with 95\% CI $[0.032, 0.076]$, similarly spanning a broad region of the SBP ranking. These mid-rank cases illustrate the value of reporting bootstrap CIs alongside point estimates: the ranking itself is accompanied by an uncertainty interval that informs how confidently one partner can be placed above or below another.

An instructive contrast emerges between R2 and R9: R2 has the second-largest age nABCD (0.061) but the smallest SBP nABCD (0.015), while R9 has the 4th-smallest age nABCD (0.017) but the largest SBP nABCD (0.110). A ranking based on either candidate effect modifier alone would reach opposite conclusions about these two regions. This asymmetry underscores the necessity of evaluating all candidate effect modifiers jointly, as visualized in Figure~\ref{fig:gusto_scatter}.

\begin{figure}[ht]
\centering
\includegraphics[width=\textwidth]{figures/fig_gusto_r8_scatter}
\caption{Joint distributional assessment for Region~8. Each point represents a partner region, positioned by nABCD for age (horizontal axis) and SBP (vertical axis). The plot presents distributional similarity as a continuum: partners closer to the lower-left corner rank low on both candidate effect modifiers (R4, R6, R13), partners further along one axis rank low on one but high on the other (R2, R9), and partners further on both axes rank high on both (R3).}
\label{fig:gusto_scatter}
\end{figure}

\subsection{Clinical Interpretation and Pooling Candidates}\label{sec:app_clinical}

For both candidate effect modifiers, $L$ is not quantified at the planning stage, so we apply the $L^*$ reverse-calculation (equation~\ref{eq:lstar}) to determine the CATE sensitivity required for the observed distributional differences to produce clinically meaningful heterogeneity ($\Delta_{\text{clin}} = 1$\%pt and $2$\%pt on the absolute 30-day mortality scale). Table~\ref{tab:gusto_lstar_joint} presents the full results for all 15 partner regions; the joint distributional pattern is visualized in Figure~\ref{fig:gusto_scatter}.

\begin{table}[ht]
\caption{Clinical interpretation ($L^*$ at $\Delta_{\text{clin}} = 1$\%pt) for all 15 partner regions, listed by region number. $L^*$ is the minimum CATE sensitivity (per year for age, per mmHg for SBP) required for the observed distributional difference to produce a treatment effect difference of $1$\%pt on the absolute 30-day mortality scale.\label{tab:gusto_lstar_joint}}
\begin{tabular*}{\textwidth}{@{\extracolsep\fill}lrrrr@{}}
\toprule
\textbf{Partner} & \textbf{nABCD\textsubscript{age}} & \textbf{nABCD\textsubscript{SBP}} & \textbf{$L^*$ age (/yr)} & \textbf{$L^*$ SBP (/mmHg)} \\
\midrule
R1  & 0.018 & 0.058 & 0.0155 & 0.0028 \\
R2  & 0.061 & 0.015 & 0.0047 & 0.0105 \\
R3  & 0.076 & 0.051 & 0.0038 & 0.0030 \\
R4  & 0.016 & 0.042 & 0.0180 & 0.0038 \\
R5  & 0.011 & 0.056 & 0.0256 & 0.0030 \\
R6  & 0.026 & 0.050 & 0.0110 & 0.0030 \\
R7  & 0.011 & 0.082 & 0.0247 & 0.0020 \\
R9  & 0.017 & 0.110 & 0.0171 & 0.0015 \\
R10 & 0.038 & 0.062 & 0.0072 & 0.0025 \\
R11 & 0.035 & 0.104 & 0.0079 & 0.0016 \\
R12 & 0.024 & 0.097 & 0.0116 & 0.0016 \\
R13 & 0.034 & 0.037 & 0.0087 & 0.0042 \\
R14 & 0.020 & 0.064 & 0.0142 & 0.0023 \\
R15 & 0.017 & 0.074 & 0.0164 & 0.0021 \\
R16 & 0.050 & 0.071 & 0.0057 & 0.0023 \\
\bottomrule
\end{tabular*}
\begin{tablenotes}
\item Regions are listed by region number; ordering carries no implicit ranking. The $L^*$ values quantify the CATE sensitivities at which these distributional differences would translate into clinically meaningful treatment effect heterogeneity.
\end{tablenotes}
\end{table}

\begin{figure}[ht]
\centering
\includegraphics[width=\textwidth]{figures/fig_gusto_r8_calibration}
\caption{Clinical interpretation ($L^*$) for Region~8. Both panels show $L^*$---the minimum CATE sensitivity that would translate the observed distributional difference into a clinically meaningful treatment effect difference---at two clinical margin values ($\Delta_{\text{clin}} = 1$\%pt, $2$\%pt). (a)~Age: $L^*$ (per year) for the partners with the largest age nABCD (R16, R2, R3). (b)~SBP: $L^*$ (per mmHg) for the partners with the largest SBP nABCD values.}
\label{fig:gusto_calibration}
\end{figure}

The joint assessment identifies three regions---R4, R6, and R13---that rank low on both candidate effect modifiers, with all six nABCD values lying in the lower portions of the observed ranges (age: 0.011--0.076; SBP: 0.015--0.110) (Table~\ref{tab:gusto_lstar_joint}; Figure~\ref{fig:gusto_scatter}, lower-left quadrant). These three regions exhibit low nABCD on both modifiers, and their required $L^*$ values (Table~\ref{tab:gusto_lstar_joint}; Figure~\ref{fig:gusto_calibration}) lie near the lower end of the range that could reasonably be considered clinically plausible for thrombolysis in AMI based on the available class evidence. Accordingly, \emph{the sponsor may reasonably prioritize R4, R6, and R13 as the leading candidates for pooling with Region~8}.

The framework does not assert binary ``poolable'' or ``not poolable'' designations. Rather, it supplies quantitative inputs---nABCD values, bootstrap confidence intervals (Table~\ref{tab:gusto_nabcd}), and $L^*$ sensitivities (Table~\ref{tab:gusto_lstar_joint})---that support the sponsor's judgment in consultation with clinical and regulatory advisors. Sponsors considering partners outside these leading candidates (R4, R6, R13) may examine the partner-specific $L^*$ values and bootstrap CI widths to evaluate the precision and plausibility of pooling on each candidate effect modifier.

